\section{\label{sec:concl}Conclusion}

We have shown how a variety of small existing and hypothetical inverse beta decay detectors can measure the incoming direction of electron-antineutrinos, given a sufficient number of events.
At these distances, one may also study the source distribution of neutrinos being produced in the reactor core. 
Furthermore, we have proposed a detector design that shows promise for substantially improving the angular resolution of these measurements within reasonable timescales for our targeted near-surface, low-standoff deployment scenario.
% include 5-20m range, ~1-5T, ~1mo.(ish)
\textbf{Specifically, our simulations indicate a potential azimuthal resolution as fine as $\bs{\sim1.7^{o}}$ with only $\bs{\sim800}$ measured \ibd\ events.}
%We have estimated that, even assuming a conservative $10\%$ detection efficiency per \ibd, a $1\mathrm{T}$-scale detector could achieve this resolution at a standoff of $10\mathrm{m}$ from a $20\mathrm{MW_{Th}}$ reactor within a live-time (detector operating time) of less than $\mathrm{48~\mathrm{hrs.}}$
We have estimated that, even assuming a conservative $10\%$ detection efficiency per \ibd, a $1-5~\mathrm{T}$-scale detector could achieve this resolution at a standoff of $5-20\mathrm{m}$ from a $20\mathrm{MW_{Th}}$ reactor within a live-time (detector operating time) on the scale of $10^1~\mathrm{hrs}$.

This is no doubt practical in proximity to a reactor (and has been accomplished by some such as PROSPECT at ORNL)~\cite{} but without very strong motivation.
% If located within $\sim~5-20$~m of the reactor, as well as observing the time evolution of burnup.
More interesting is the case for finding unknown reactors at ranges of $10-100$~km.
The 10-km scale will require fiducial volumes of order several tens of tons, and 100~km will require a kiloton fiducial volume.
Combining the direction-finding and range-finding, we see that at least in principle, one may achieve blind recognition of a reactor's existence, azimuth, distance and power~\cite{}. % reactor-ranging reference steve dye and bob svoboda john paper
While not an easy task, it is a remarkable capability unique to neutrinos \dash the only stable elementary particles which ``know'' their point of origin!

~